\section{A fixed-level high-field application}\label{sec:fixed}

\subsection{Redundancy eight}

\begin{proposition}[Generic three-marker bottom cover]\label{prop:r8-bottom}
After three contractions, the generic bottom object is the R5
off-diagonal $(2,2)$ curve.  The diagonal and branch deletion costs
twelve and three marked roots cost eighteen, so
$\delta=30$.  The inequality
\[
       q+1-2\sqrt q>30
\]
holds at every prime power $q\geq43$.  The first-polar line has
transverse/collision budget
\[
       3+1+10=14,
\]
where ten is the ramification degree of a moving $g^4_6$.
\end{proposition}

\begin{proof}
Markers do not change the bottom curve or its geometric monodromy;
each removes at most the six ordered pairs in its unique cubic fiber.
The point bound is the \(\kappa=1\) Aubry--Perret form used throughout.
A noncontained line meets the lower
rank-two catalecticant carrier in degree at most three, and the lower
binary central point in at most one parameter.  After fixed factors are
removed, self-collision is ramification of a separable $g^4_6$, of
degree $(4+1)(6-4)=10$.  Separability follows because the
$p$-th-power subspace has dimension smaller than the five-dimensional
series.
\end{proof}

\begin{definition}[Three-marker lower-package gate]
\label{def:r8-lower-package}
The assertion $\mathrm{LP}(6,1)$ says that the pointed R7 splitting
incidence, after the persistent rank-two carrier, the binary central
point, and the marked-collision divisor are removed, has the
geometrically integral identity twists and deletion bounds used in
Proposition~\ref{prop:r8-bottom}.  It includes equations and
monodromy checks for the recursive cyclic, wild, inseparable, gcd, and
branch strata; a finite-field regression is not a substitute for those
geometric data.
\end{definition}

The ordered cover retains the marker labels needed for monodromy, but
the contraction itself descends to the unordered marker cubic.  Indeed,
for affine markers \(r,s,t\),
\begin{equation}\label{eq:r8-marker-cubic}
 a_i^{\langle r,s,t\rangle}
 =a_{i+3}-(r+s+t)a_{i+2}
 +(rs+rt+st)a_{i+1}-rst\,a_i .
\end{equation}
Thus it depends only on the elementary symmetric functions of the
markers, equivalently on the monic cubic
\(T^3-(r+s+t)T^2+(rs+rt+st)T-rst\).  This separates the symmetric
contraction parameter from the ordered-root cover used below.

\begin{lemma}[Bottom strata]\label{lem:r8-bottom-strata}
For the two-step contraction of a pointed sextic, the lower rank,
gcd-one, cyclic, wild, inseparable, and branch strata are exactly the
closed schemes displayed below.  The reversed homogeneous chart
introduces no additional stratum.
\end{lemma}

\begin{proof}
For a pointed sextic
$f=(a_0,\ldots,a_6)$ with old marker $x$, contract successively:
\[
 b_i(r)=a_{i+1}-ra_i,\qquad
 c_i(r,s)=b_{i+1}(r)-sb_i(r).
\]
For $c=(c_0,\ldots,c_4)$ put
\[
 H_c=\begin{pmatrix}c_0&c_1&c_2&c_3\\
                     c_1&c_2&c_3&c_4\end{pmatrix}.
\]
On the rank-two-kernel chart, if $p_0,p_1$ span $\ker H_c$, the bottom
ordered-root curve is
\begin{equation}\label{eq:r8-bottom-cover}
 G_c(u,v)=
 \frac{p_0(u)p_1(v)-p_1(u)p_0(v)}{[u,v]}=0.
\end{equation}
The rank-at-most-two secant carrier, including the gcd-two and
rank-drop boundary strata, is
\begin{equation}\label{eq:r8-secant-carrier}
 D(c)=\det\begin{pmatrix}
 c_0&c_1&c_2\\c_1&c_2&c_3\\c_2&c_3&c_4
 \end{pmatrix}=0.
\end{equation}
For an exact common root $z$, the gcd-one incidence is cut out by
\begin{equation}\label{eq:r8-gcd-one}
 \operatorname{rank}
 \begin{pmatrix}H_c\\1&z&z^2&z^3\end{pmatrix}\leq2.
\end{equation}
with the evident homogeneous row at infinity.  After substitution,
\eqref{eq:r8-secant-carrier} has degree at most three in $s$, while
each $3$-minor in \eqref{eq:r8-gcd-one} has degree at most two for
fixed $z$.

The remaining reducible-monodromy carriers are explicit.  Put
$z_i=\binom4i c_i$.  Outside characteristics two and three the cyclic
closure is
\begin{equation}\label{eq:r8-cyclic-veronese}
 [u:v:w]\longmapsto
 [u^2:2uv:v^2+2uw:2vw:w^2],
\end{equation}
a degree-four projected Veronese surface containing no line.  In
characteristic three it is the degree-four cone
$\Join(e_2,C_4)$; its only lines are rulings, and the consecutive-row
overlap excludes every ruling as a polar line.  In characteristic two
it is
\begin{equation}\label{eq:r8-cyclic-binary}
 \Pi_{\rm cyc}=\PP\langle e_1,e_2,e_3\rangle,\qquad c_0=c_4=0.
\end{equation}
whose unique contained polar-line component is
$\cN_3=\PP\langle e_2,e_3\rangle$.  The inseparable degeneration is
the derivative-minor locus of $p_0,p_1$ and introduces no further
stratum.  Outside characteristics two and three a nonconstant
degree-three map cannot factor through Frobenius.  In characteristic
three an inseparable cubic pencil is projectively
$\langle X^3,Y^3\rangle$; the two Hankel rows then force
$c_0=c_1=c_3=c_4=0$, the vertex $e_2$ of the wild cone.  In
characteristic two, after cancelling a nonzero pencil gcd, an inseparable
ratio has degree at most two and is a rational function of $t^2$.
Homogenizing its numerator and denominator to cubics supplies the
same linear factor to both, so it belongs to the exact-gcd-one
boundary.  Finally, the branch equation on this chart is
$p_0'p_1-p_0p_1'=0$.  These carrier and branch loci are the
scheme-theoretic images of the displayed incidences (equivalently their
Fitting-minor ideals); the reversed homogeneous chart supplies infinity.
Thus these are closed strata over the coefficient ring, not a
classification only of their rational points.
\end{proof}

\begin{lemma}[Bottom monodromy and deletion]\label{lem:r8-monodromy}
Off the strata of Lemma~\ref{lem:r8-bottom-strata}, the bottom
ordered-root curve is geometrically integral of arithmetic genus one
and has deletion degree at most $30$.  On the exact-gcd-one stratum,
the residual deck graph supplies a squarefree pair avoiding the three
markers for $q\geq43$.
\end{lemma}

\begin{proof}
Off \eqref{eq:r8-secant-carrier}--\eqref{eq:r8-cyclic-binary}, the
separable cubic map is noncyclic, so its geometric
monodromy is $S_3$.  Transitivity on ordered pairs of distinct roots
makes \eqref{eq:r8-bottom-cover} geometrically integral.  It has
arithmetic genus one.
The diagonal costs four, the degree-four different costs at most
eight off-diagonal ordered pairs, and each of the three marker fibers
$x,r,s$ costs at most six.  Thus
\begin{equation}\label{eq:r8-deletion}
 \delta\leq4+8+3\cdot6=30.
\end{equation}
Hasse--Weil supplies a surviving rational point for every $q\geq43$.

It remains to show that the strata removed above do not hide another
component.  On an exact-gcd-one bottom stratum write the cubic pencil
as $\ell_zQ$, with $Q$ a basepoint-free pencil of quadratics.  In the
separable case its deck involution $\iota$ is rational, so its graph
has $q+1$ rational points.  At most two are fixed and at most two
ordered graph points involve each of $z,x,r,s$.  A point remains for
$q\geq43$, giving three distinct roots outside the markers.  The
inseparable degree-two case is on the rank/boundary carrier.
\end{proof}

\begin{lemma}[Two-marker selection]\label{lem:r8-two-marker}
Fix two old markers $x,r$.  Outside the declared contained strata,
the bad divisor on the internal $s$-line has degree at most $19$.
For $q\geq43$ one can choose $s$ so that the lower package supplies a
split squarefree quartic avoiding $x$ and $r$.
\end{lemma}

\begin{proof}
The complete divisor table on the $s$-line is
\[
\begin{array}{c|c|c}
\text{failure}&\text{defining mechanism}&\text{degree}\\\hline
\text{rank at most two}&D(c(s))&3\\
\text{cyclic/wild monodromy}&\text{pullback of
\eqref{eq:r8-cyclic-veronese}--\eqref{eq:r8-cyclic-binary}}&4\\
\text{new gcd root }s&\operatorname{Ram}(g^2_d),\ d\le4&6\\
\text{gcd root }x\text{ or }r&\text{one nonzero minor in
\eqref{eq:r8-gcd-one}}&2+2\\
\text{marker collision}&s=x,\ s=r&2
\end{array}
\]
In characteristic two the cyclic entry has degree one, so degree four
is a uniform bound.  The total is
\begin{equation}\label{eq:r8-inner-budget}
 3+4+6+2+2+2=19<q+1.
\end{equation}
The rank pullback cannot be the whole line:
Proposition~\ref{prop:r6-secant} identifies that containment exactly
with quadratic gcd of the original net.  The only possible contained
cyclic line is the declared binary nucleus.  For a fixed marker, not
all minors in \eqref{eq:r8-gcd-one} can vanish identically without
giving the original
quartic net that fixed factor: otherwise the hyperplanes
$W\cap\ker(\operatorname{ev}_s)$ and
$W\cap\ker(\operatorname{ev}_x)$ agree for dense $s$, making
the evaluation map constant.  Choose one nonzero quadratic minor.
For self-collision, take a generic pencil projection of the moving
$g^2_d$, $d\leq4$.  It is separable: in characteristic two the only
dimension-three equality case is the pure-square space
$\langle1,t^2,t^4\rangle$, which the consecutive Hankel overlap
equations do not realize; all other $p$-power spaces have smaller
dimension.  Riemann--Hurwitz gives degree at most $2d-2\leq6$, and
every self-collision maps into this ramification divisor.
Choosing $s$ outside this divisor and applying the preceding
gcd-zero or gcd-one bottom analysis gives a split squarefree quartic
avoiding the two old markers.  This proves the two-old-marker lemma.
\end{proof}

\begin{lemma}[Outer marker selection]\label{lem:r8-outer-selection}
For a pointed sextic outside the persistent and modular contained
strata, the bad divisor on the outer first-marker line has degree at
most $15$.  A surviving outer marker lifts the witness of
Lemma~\ref{lem:r8-two-marker} to a split squarefree quintic avoiding
the original marker.
\end{lemma}

\begin{proof}
The lower rank-two pullback costs three
parameters and the binary nucleus costs one.  If a contracted net has
exact gcd $\ell_z$, the condition $z=r$ maps into the ramification of
a generic separable pencil projection of the moving $g^3_d$, $d\leq5$.
Riemann--Hurwitz bounds it by $2d-2\leq8$.  The condition $z=x$ is
again \eqref{eq:r8-gcd-one}, now with the quartic evaluation row, and
costs at most two
because not all of its quadratic minors vanish off the fixed-marker
boundary, by the same evaluation-hyperplane argument.  Together
with $r=x$, the internal divisor has degree
\begin{equation}\label{eq:r8-outer-budget}
 3+1+8+2+1=15.
\end{equation}
If the contracted net has trivial gcd, the two-old-marker lemma applies.
If it has exact gcd with $z\notin\{x,r\}$, the ordered-pair argument has three
$(1,1)$ and two $(2,1)/(1,2)$ bad equations, at most
$12(q+1)<(q-2)(q-3)$.  The cases $z=x,r$ are precisely the excluded
quadratic and ramification divisors.
Lifting first by $\ell_s$ and then by $\ell_r$ gives a split
squarefree quintic avoiding $x$.
\end{proof}

\begin{proposition}[Pointed lower package]\label{prop:r8-lp61}
The assertion $\mathrm{LP}(6,1)$ holds.  On the outer first-marker
line its exceptional divisor has degree at most fifteen; on the
recursive second-marker line the corresponding divisor has degree at
most nineteen.  Its only positive-genus identity twist is the
geometrically integral off-diagonal $(2,2)$ curve of genus at most one
with deletion degree $30$.
\end{proposition}

\begin{proof}
Lemmas~\ref{lem:r8-bottom-strata} and \ref{lem:r8-monodromy} identify
the bottom strata and the unique positive-genus identity twist.
Lemmas~\ref{lem:r8-two-marker} and \ref{lem:r8-outer-selection} give
the two parameter-line budgets and the coherent squarefree lift.  For
auditability, the complete exhaustion is summarized here.

\begin{center}
\scriptsize
\setlength{\tabcolsep}{3pt}
\begin{tabular}{p{0.12\textwidth}p{0.21\textwidth}p{0.08\textwidth}
                p{0.19\textwidth}p{0.21\textwidth}}
\toprule
Stratum & Defining equation & Degree & Contained case &
Transverse conclusion\\
\midrule
rank at most two & \eqref{eq:r8-secant-carrier} &
$3$ in the marker & quadratic-gcd persistent carrier &
finite intersection of degree at most $3$\\
cyclic or wild & \eqref{eq:r8-cyclic-veronese} or
\eqref{eq:r8-cyclic-binary} & at most $4$ &
only the binary nucleus line & finite intersection of degree at most
$4$\\
exact gcd one & minors in \eqref{eq:r8-gcd-one} &
$2$ per fixed marker & fixed-factor boundary &
the residual deck graph supplies a squarefree pair\\
inseparable & derivative minors of $p_0,p_1$ & absorbed above &
wild vertex, rank drop, or binary nucleus & no additional stratum\\
geometric $S_3$ & \eqref{eq:r8-bottom-cover} &
bidegree $(2,2)$ & only the preceding carriers &
integral genus-one cover with deletion degree $30$\\
marker collision & diagonals and ramification &
$1$ per diagonal; at most $6$ for self-collision &
declared collision boundary & finite divisor on the marker line\\
\bottomrule
\end{tabular}
\end{center}

The gcd has degree zero, one, or at least two, and a separable
gcd-zero cubic cover has monodromy $S_3$ or $C_3$; hence the displayed
list is exhaustive.  No certificate or finite-field regression is
used for this integrality assertion.  Notice also that the outer and
recursive parameter budgets $15$ and $19$ are already strictly below
$q+1$ for $q\geq19$; the threshold $43$ is imposed solely by the
genus-one deletion inequality~\eqref{eq:r8-deletion}, not by a hidden
parameter-choice
bound.
\end{proof}

\begin{proposition}[Modular lifts at degree seven]\label{prop:r8-modular}
The consecutive-row Lucas calculation leaves no nonzero lift in
characteristic two, the line $\PP\langle e_2,e_5\rangle$ in
characteristic three, and the line
$\PP\langle e_3,e_4\rangle$ in characteristic five.  The latter two
lines are shallow in the range $q\geq43$.
\end{proposition}

\begin{proof}
Lucas' criterion applied to the consecutive binomial rows gives the
three stated supports; the exact row reductions are replayed by
Certificate R8.  In characteristic five, the homogenization of
$t^5-t$ supplies six distinct projective roots and satisfies the two
required middle-coefficient equations throughout the lift.

In characteristic three, $\PGL_2(q)$ is transitive on the rational
points of $\PP\langle e_2,e_5\rangle$, so it suffices to treat $e_2$.
Choose nonzero $a,b,c$ with no two equal up to sign and
$a^2+b^2+c^2=0$.  The nonsingular conic has $q+1$ points; deleting the
three coordinate and six equality/antiequality sections costs at most
eighteen, so a choice exists for $q\geq27$.  The reciprocal pairs
$\pm a^{-1},\pm b^{-1},\pm c^{-1}$ give six distinct roots and the two
required elementary-symmetric coefficients vanish.
\end{proof}

\begin{proposition}[Degree-seven contained components]
\label{prop:r8-contained}
The assertion $\mathrm{CC}(7,1)$ for the structural components of
Definition~\ref{def:r8-lower-package} holds in every characteristic.
The rank-two component lifts to the upper rank-two carrier; the lower
binary central point has no nonzero coherent lift; and the
marked-collision scheme is never the whole polar line.  The only
additional nucleus lifts are
\[
 \PP\langle e_2,e_5\rangle\quad\text{in characteristic three},
 \qquad
 \PP\langle e_3,e_4\rangle\quad\text{in characteristic five},
\]
and both are shallow in the range $q\geq43$.
\end{proposition}

\begin{proof}
The rank-two assertion is
Proposition~\ref{prop:contained-rank-two} with $n=7$.  For the binary
central point $e_3$, requiring both consecutive septic rows to be
supported at its single coordinate permits $a_3$ in the first row and
$a_4$ in the second; their overlap is zero.  Thus its projective
coherent lift is empty.  The full nucleus-support overlap is exactly
Proposition~\ref{prop:r8-modular}, which also supplies the shallow
witnesses in characteristics three and five.

It remains to exclude an identically colliding series.  Remove the
fixed divisor from the five-dimensional $W_f$ and let the moving
degree be $d\leq6$.  If its projective map were inseparable in
characteristic $p$, it would factor through Frobenius, so its section
space would have dimension at most
\[
 \left\lfloor d/p\right\rfloor+1\leq4,
\]
contrary to $\dim W_f=5$.  A generic pencil projection is therefore a
separable map of degree at most six.  Every marked self-collision is a
ramification point of that projection, so Riemann--Hurwitz bounds the
collision divisor by $2d-2\leq10$.  It is finite and cannot be an
additional contained component.
\end{proof}

\begin{remark}[Exact remaining R8 gate]\label{rem:r8-contained}
Proposition~\ref{prop:r8-contained} closes the contained-component
problem, while Proposition~\ref{prop:r8-lp61} closes the pointed lower
package.  The certificate checks the nucleus overlaps, witnesses, and
numerical budgets; the geometric integrality and recursive carrier
equations are proved separately above.
\end{remark}

\begin{theorem}[Redundancy eight]\label{thm:r8}
For every prime power $q\geq43$, the deep
syndromes of $\PRS(q-7)$ are
exactly the persistent tangent and conjugate-secant families, of total
size $q(q+1)^2/2$.  Their orbit law is $T/T^7$ modulo inversion and
Frobenius.  The tangent family splits exactly in characteristic seven.
\end{theorem}

\begin{proof}
Outside the contained locus, Proposition~\ref{prop:r8-bottom} chooses a
good first contraction, and Proposition~\ref{prop:r8-lp61} supplies a
three-marker witness.  Propositions~\ref{prop:r8-modular} and
\ref{prop:r8-contained} eliminate every nonpersistent contained lift.
The imported covering-radius theorem promotes the resulting split-free
classification to deep holes.
\end{proof}

This theorem makes no assertion about additional deep orbits for
$q<43$.

\subsection{Redundancy nine}

\begin{definition}[R9 pointed domain]\label{def:r9-pointed-domain}
The assertion $\mathrm{LP}(7,1)$ is made on the pointed splitting
incidence after removing the persistent rank-two carrier, the complete
characteristic-specific Lucas-nucleus support, and the fixed-marker and
self-collision components.  At every recursive level, a polar line
contained in one of these supports belongs to that declared closed
boundary; the parameter-divisor argument applies only to the
noncontained line.  Thus a section that vanishes identically is never
counted as a finite divisor of its nominal degree.
\end{definition}

\begin{proposition}[Four-marker pointed lower package]
\label{prop:r9-budgets}
The assertion $\mathrm{LP}(7,1)$ holds for every $q\geq53$.
Its bottom identity twist is the geometrically integral off-diagonal
$(2,2)$ curve of Proposition~\ref{prop:r8-lp61}, now with four
forbidden markers and deletion degree
\begin{equation}\label{eq:r9-deletion}
       4+8+4\cdot6=36.
\end{equation}
The three recursive parameter-line budgets, from the bottom upward,
are at most $22$, $18$, and $18$.  At the top degree-eight
first-polar line the transverse/component budget is at most
\begin{equation}\label{eq:r9-top-budget}
       3+2+12=17.
\end{equation}
No geometric-integrality, gcd, or contained-component assertion is
included merely as a number in these budgets.
\end{proposition}

\begin{proof}
We retain the notation and the scheme-theoretic strata
\eqref{eq:r8-bottom-cover}--\eqref{eq:r8-cyclic-binary}.  Their
homogeneous incidence ideals commute with further
contraction, so adding a marker changes only the corresponding
evaluation and diagonal pullbacks; it does not change the bottom
cover.  We prove the package in three steps.

\emph{Three-pointed quartic net.}
Let $W$ be a trivial-gcd Hankel net of quartics, let
$x_1,x_2,x_3$ be distinct prescribed markers, and exclude the
characteristic-two nucleus already listed in
\eqref{eq:r8-cyclic-binary}.  Contract at a
fourth parameter $s$.  On the resulting cubic pencil the rank
pullback \eqref{eq:r8-secant-carrier}, the cyclic/wild pullback
\eqref{eq:r8-cyclic-veronese}--\eqref{eq:r8-cyclic-binary}, and
self-collision
have degrees at most $3$, $4$, and $6$, respectively.  For each
$x_i$, choose a nonzero quadratic minor in \eqref{eq:r8-gcd-one}; the
three fixed-root
pullbacks cost at most $6$, and the three equations $s=x_i$ cost $3$.
Thus the parameter divisor has degree at most
\begin{equation}\label{eq:r9-three-point-budget}
             3+4+6+3\cdot2+3=22.
\end{equation}
None of these equations vanishes identically.  For the rank and
cyclic carriers this is the no-contained-line calculation following
\eqref{eq:r8-inner-budget}.  For a fixed marker, identical vanishing
of all the minors in \eqref{eq:r8-gcd-one} would make the kernels of
$\operatorname{ev}_{s}|_W$ and $\operatorname{ev}_{x_i}|_W$ agree
for dense $s$.  The two nonzero functionals would then be
proportional, making the projective evaluation map of the
three-dimensional trivial-gcd net $W$ constant, a contradiction.
For self-collision, a generic pencil projection of the moving
$g^2_d$, $d\leq4$, is separable by the pure-power dimension argument
following \eqref{eq:r8-inner-budget}, and Riemann--Hurwitz gives
degree at most $6$.

Away from this divisor, a gcd-zero cubic pencil is on the geometric
$S_3$ stratum.  Its ordered-pair cover \eqref{eq:r8-bottom-cover} is
geometrically integral
of genus at most one.  The diagonal and different delete at most
$4+8$ points, and each of $x_1,x_2,x_3,s$ deletes at most the six
ordered pairs in its cubic fiber.  This proves
\eqref{eq:r9-deletion}.  If instead the
cubic pencil has exact gcd $\ell_z$, cancel it and use the graph of
the rational deck involution of the residual quadratic pencil.  The
graph has $q+1$ rational points; at most two are fixed, and at most two
involve each of $z,x_1,x_2,x_3,s$.  Since $q+1>12$, a graph point
remains.  In the inseparable quadratic case the characteristic is two
and, after a root-coordinate change, the residual pencil is
$\langle T^2,U^2\rangle$.  On the affine chart
$\ell_z=t-z$, its two kernel cubics have coefficient vectors
\[
             (-z,1,0,0),\qquad(0,0,-z,1).
\]
The two overlapping Hankel rows give successively
\[
 c_1=zc_0,\quad c_2=z^2c_0,\quad
 c_3=z^3c_0,\quad c_4=z^4c_0.
\]
Hence the determinant \eqref{eq:r8-secant-carrier} vanishes; the
homogeneous $z=\infty$
chart gives the endpoint.  Thus the inseparable case really lies on
the rank-drop carrier.  Multiplication by $\ell_s$ gives a split
squarefree quartic avoiding $x_1,x_2,x_3$.

\emph{Two-pointed quintic series.}
Let $f\in S_6$ carry two distinct markers $x_1,x_2$ and lie outside
the lower persistent, binary-nucleus, fixed-marker, and collision
boundary of Definition~\ref{def:r9-pointed-domain}; contract at $r$.
The lower rank-two and binary-nucleus pullbacks then cost at most
$3+1$.  If the contracted quartic net has exact gcd $\ell_z$, the
condition $z=r$ lies in the ramification divisor of a generic
separable projection of the moving $g^3_d$, $d\leq5$, and costs at
most $8$.  The conditions $z=x_i$ are the two fixed-evaluation
incidences \eqref{eq:r8-gcd-one}, each of degree at most two;
$r=x_i$ contributes two
diagonals.  Hence
\begin{equation}\label{eq:r9-two-point-budget}
             3+1+8+2\cdot2+2=18.
\end{equation}
The same rank, evaluation-hyperplane, and separable-projection
arguments show that no listed pullback is the whole $r$-line off the
declared boundary.

For a trivial-gcd contracted net, the preceding three-pointed lemma
applies to $x_1,x_2,r$.  For exact gcd $\ell_z$ away from these
markers, cancel $\ell_z$.  In the ordered-pair parametrization of the
residual cubic hyperplane, avoiding $z,x_1,x_2,r$ gives four
bidegree-$(1,1)$ equations; the two root self-collisions have
bidegrees $(2,1)$ and $(1,2)$.  Their rational zero sets contain at
most $14(q+1)$ ordered pairs, whereas
\[
             (q-3)(q-4)>14(q+1)\qquad(q\geq53).
\]
Thus a split squarefree quartic remains.  The cases
$z\in\{x_1,x_2,r\}$ are exactly the excluded fixed-evaluation and
ramification divisors.  Lifting by $\ell_r$ gives a split squarefree
quintic avoiding $x_1,x_2$.

\emph{One-pointed sextic series.}
Now let $f\in S_7$ have marker $x$ and lie in the open pointed domain
of Definition~\ref{def:r9-pointed-domain}; contract at $r$.  The
pullback of the lower rank-two carrier has degree at most three.  The
complete consecutive-row Lucas calculation gives at most two linear
nucleus-support pullbacks.  A whole binary-central line belongs to the
binary Lucas-nucleus component, and the other whole-line possibilities
are exactly the remaining declared modular supports, so these
pullbacks are finite on the pointed domain.  After removing fixed
factors, a generic
projection of the moving $g^4_d$, $d\leq6$, is separable: an
inseparable series of degree at most six has dimension at most
$\lfloor6/p\rfloor+1<5$.  Its ramification degree is at most
$2d-2\leq10$, and it contains the collision in which the new gcd root
equals $r$.  The fixed-root incidence at $x$ costs at most two, and
$r=x$ costs one.  The total is
\begin{equation}\label{eq:r9-one-point-budget}
             3+2+10+2+1=18.
\end{equation}
Choose $r$ off this divisor and apply the two-pointed quintic result
with markers $x,r$.  Its quintic avoids both, so multiplication by
$\ell_r$ gives a split squarefree sextic in $W_f$ avoiding $x$.
This is $\mathrm{LP}(7,1)$.

Finally, let $f\in S_8$ lie outside the persistent, modular, and
self-collision components of $\mathrm{CC}(8,1)$.  Its first-polar line
meets the lower
rank-two carrier in degree at most three and the complete linear
nucleus support in degree at most two.  The new-marker collision is
ramification of a separable $g^5_d$, $d\leq7$: an inseparable series
would have dimension at most $\lfloor7/p\rfloor+1<6$, and
Riemann--Hurwitz gives at most $2d-2\leq12$.  This proves
\eqref{eq:r9-top-budget} and
excludes an additional contained collision component.

The gcd at every recursive level has degree zero, one, or at least
two.  The last case lies in the degree-appropriate rank-two
catalecticant carrier, whose bottom equation is
\eqref{eq:r8-secant-carrier} and whose
contained-line pullback is
Proposition~\ref{prop:contained-rank-two}; the degree-one case is
treated after cancellation.  A separable gcd-zero cubic map has geometric
monodromy $S_3$ or $C_3$, with the latter and every inseparable case
in \eqref{eq:r8-cyclic-veronese}--\eqref{eq:r8-cyclic-binary} or the
derivative-minor boundary.  Hence the strata are
exhaustive.  Since
\[
             q+1-2\sqrt q>36
\]
at $q=53$ and the left side is increasing for $q>1$,
\eqref{eq:r9-deletion} leaves a
rational point for every $q\geq53$.  The parameter budgets
\eqref{eq:r9-three-point-budget}--\eqref{eq:r9-one-point-budget} and
\eqref{eq:r9-top-budget} are all smaller than $q+1$.  The first prime
power satisfying the genus-one inequality is $53$.
\end{proof}

The only substantial modular carrier occurs in characteristic seven
and is a projective binary-quartic space.  Fixing a split quintic
$P(t)=\sum_{i=0}^5p_it^i$ and a carrier point
$h=(a_0,\ldots,a_4)$, put
\[
 H_j=\sum_{i=0}^4a_ip_{i+j}
\]
with $p_k=0$ outside $0\leq k\leq5$, and
\[
 D=H_0H_2-H_1^2,\quad
 N_s=H_{-1}H_2-H_0H_1,\quad
 N_u=H_{-1}H_1-H_0^2.
\]
Off $D=0$, the unique residual quadratic is
\[
       Dt^2-N_st+N_u,
\]
with branch polynomial $K=N_s^2-4N_uD$.  The exact residual identities
are proved next; the geometric exhaustion of transported slices is
kept as a separate printed gate.

\begin{proposition}[Residual quadratic identities]\label{prop:r9-residual}
With $p_k=0$ outside $0\leq k\leq5$, the two consecutive Hankel
equations for
\[
       P(t)(t^2-st+u)
\]
are
\[
H_0-sH_1+uH_2=0,\qquad
H_{-1}-sH_0+uH_1=0.
\]
On $D\ne0$, Cramer's rule gives
\[
       s=N_s/D,\qquad u=N_u/D.
\]
The branch divisor is $K=0$, fixed/residual collision is
$\operatorname{Res}(P,Dt^2-N_st+N_u)=0$, and fixed-root collision is
$\operatorname{Disc}(P)=0$.
\end{proposition}

\begin{proof}
Expand the two consecutive coefficients annihilated by the carrier
quartic.  The displayed linear system has determinant $D$, so Cramer's
rule gives the residual factor.  Its discriminant is $K$.  The two
resultants have their usual intrinsic collision meanings; consequently
the construction commutes with scalar, $\PGL_2$, and Frobenius
transport.
\end{proof}

\begin{proposition}[One-moving-root slice]\label{prop:r9-slice}
Let $P=R(t)(t-x)$ with four fixed distinct roots.  Then
$D,N_s,N_u$ have degree at most two in $x$, and $K$ has degree at most
four.  If the squarefree reduction of $K$ is nonconstant and not a
square, the normalization of $y^2=K(x)$ is geometrically integral of
genus at most one.
\end{proposition}

\begin{proof}
Each $H_j$ is affine in $x$, so its $2\times2$ minors are quadratic and
$K$ is quartic.  In odd characteristic the squarefree nonsquare
equation defines a geometrically integral double cover; the
hyperelliptic genus formula gives
$\lfloor(\deg K_{\mathrm{red}}-1)/2\rfloor\leq1$.
\end{proof}

\begin{proposition}[Binary-quartic slices and rational base selection]
\label{prop:r9-components}
Every nonzero geometric binary quartic has a four-root slice whose
reduced residual cover is geometrically integral of genus at most one.
If the quartic is rational over a characteristic-seven field
$\F_q$ with $q>102$, the four roots may be chosen distinct and
rational.
\end{proposition}

\begin{proof}
Let
\[
 \mathcal H=\PP(\Sym^4E),\qquad
 \mathcal B=\operatorname{Conf}_4(\PP^1).
\]
Over $\mathcal H\times\mathcal B$, the residual formulas define a
universal branch quartic $K_{h,\mathbf r}(x)$.  Its subdiscriminants
cut out the locus where the squarefree part drops degree.  Hence the
good-slice incidence
\[
 \mathcal U=\{(h,\mathbf r):K_{h,\mathbf r,\mathrm{red}}
 \text{ is nonconstant and separable}\}
\]
is a $\PGL_2$-stable open set.  The geometric component problem is the
surjectivity of $\mathcal U\to\mathcal H$.  One successful
specialization would show only generic nonemptiness; fiberwise
surjectivity requires an open cover of quartic moduli.

\begin{figure}[H]
\centering
\[
\begin{array}{ccccc}
h_L
&\xrightarrow{\ \text{six sections}\ }&
(\Delta_1,\ldots,\Delta_6)
&\xrightarrow{\ \sum b_i\Delta_i=1\ }&
\displaystyle\bigcup_iD(\Delta_i)=\mathbb A^1_L\\[6pt]
\{4,\ 3+1,\ 2+2,\ 2+1+1\}
&\xrightarrow{\ \operatorname{disc}=3,3,1,5\ }&
\multicolumn{3}{c}{\text{good slices on every boundary stratum}.}
\end{array}
\]
\caption{The redundancy-nine good-slice argument.  The B\'ezout
identity covers the entire squarefree normal-form atlas; four
root-partition calculations cover its boundary.  The conclusion is
fiberwise surjectivity, not generic success of one specialization.}
\label{fig:r9-slice-cover}
\end{figure}

The quotient has one one-dimensional squarefree stratum and four
multiple-root boundary strata.  On the standard even normal-form atlas
for the squarefree stratum, write $h_L=[1,0,L,0,1]$.  The six
four-point sets and the coefficient vectors of $\Delta_i(L)$ and
$b_i(L)$ are printed in Appendix~\ref{app:r9-slices}.  They are six
test sections of the universal base family over this atlas.  Pullback
of the branch discriminant gives $\Delta_i(L)$, and direct
multiplication gives
\begin{equation}\label{eq:r9-bezout}
             \sum_{i=1}^6 b_i(L)\Delta_i(L)=1
             \quad\text{in }\F_7[L].
\end{equation}
Thus the principal opens $D(\Delta_i)$ cover the normal-form atlas;
these are six sections, not six isolated successful tests.  Transport
moves the quartic and its four-point base together, so the existence
statement descends despite the many-to-one normal form.  On the
boundary, root multiplicity gives exactly the partitions
$4,3+1,2+2,2+1+1$: there are at most three support points, and triple
transitivity gives one orbit per partition.  Their reduced branch
discriminants are $3,3,1,5$, respectively.  Thus every boundary orbit
also meets $\mathcal U$, proving surjectivity.

Now fix $h\in\mathcal H(\F_q)$.  Over
$\F_q(\mathbf r)$ let $d(h)\geq1$ be the number of distinct roots of
the generic $K_{h,\mathbf r}$.  Define $B_h(\mathbf r)$ to be its
$d(h)$-th subdiscriminant, equivalently the corresponding nonzero
Hermite principal minor.  This is a polynomial, not a discriminant
after an unspecified gcd division.  Its nonvanishing after
specialization guarantees a nonconstant separable squarefree part.
Surjectivity over the algebraic closure proves $B_h\ne0$.
The affine chart of four finite roots is dense in $\mathcal B$, so this
nonempty good open meets the chart used below.
The bad affine base locus is contained in
\begin{equation}\label{eq:r9-base-bad}
 B_h(\mathbf r)\prod_{i<j}(r_i-r_j)=0,
\end{equation}
so rational selection is avoidance of a proper divisor in configuration
space, not a second monodromy assertion.

Every coefficient of $R_{\mathbf r}(t)(t-x)$ is affine in $x$, of
degree at most one in each $r_i$ and total $\mathbf r$-degree at most
four.  Hence $D,N_s,N_u$ have individual degree at most two and total
degree at most eight; $K$ has individual degree at most four and total
degree at most sixteen.  Every quartic subdiscriminant is a Hermite
minor of coefficient-degree at most six.  Therefore
\begin{equation}\label{eq:r9-base-degree}
 \deg_{r_i}B_h\leq24,\qquad \deg B_h\leq96.
\end{equation}
Multiplying by the four-root Vandermonde gives a nonzero polynomial of
total degree at most $102$.  Schwartz--Zippel leaves a rational
ordered base with distinct roots whenever $q>102$.
\end{proof}

\begin{proposition}[Good-slice deletion]\label{prop:r9-deletion}
Once a rational base with an integral slice is fixed, the moving/fixed
diagonal, determinant, branch, four fixed-root collisions, and
moving-root collision delete at most, respectively,
\[
       8,\ 4,\ 4,\ 8,\ 8
\]
points on the normalized double cover.  Their total is $32$, so
$q+1-2\sqrt q>32$ supplies a split witness.
\end{proposition}

\begin{proof}
The projection of the normalized residual cover to the moving-root line
has degree two.  Therefore the degree-four fixed-root diagonal and the
degree-two determinant pull back to at most eight and four points.  The
degree-four branch divisor is ramified and contributes at most four
points.  For each of the four fixed roots the residual-root collision is
linear in the moving parameter; its pullback contributes at most two
points, for a total of eight.  The moving-root resultant has degree at
most four and contributes at most eight points after pullback.  Removing
their union from the normalized genus-at-most-one slice proves the claim.
\end{proof}

\begin{proposition}[Characteristic-seven finite bridges]
\label{prop:r9-char7}
At $q=7$, the formal carrier has exactly $819$ projective rootless
quartics.  At $q=49$, Certificate R9-49 exhausts all
\[
       \frac{49^5-1}{48}=5{,}884{,}901
\]
projective carrier points and finds a split witness for each.
\end{proposition}

\begin{proof}
At $q=7$, a split squarefree degree-seven projective divisor is the
complement of one of the eight rational points.  Inclusion--exclusion
gives
\[
\frac{7^5-1}{6}-8\frac{7^4-1}{6}
+28\frac{7^3-1}{6}-56\frac{7^2-1}{6}
+70\frac{7-1}{6}=819.
\]
The $q=49$ clause is deliberately certificate-only: one lexicographic
five-root search stops after every carrier point has a witness.  The
residual algebra is independently replayed, but there is no second
exhaustive carrier implementation.
\end{proof}

\begin{proposition}[R9 persistent orbit law]\label{prop:r9-orbits}
Put $d=\gcd(8,q+1)$.  The sigma-fiber projective orbits are the
inversion orbits in $C_d=T/T^8$, and coefficient Frobenius acts by
multiplication by the characteristic.  The tangent fiber is transitive
unless the characteristic is two, when its zero and nonzero parts are
separate.
\end{proposition}

\begin{proof}
The nonsplit torus acts by eighth powers and its normalizer inverts the
torus coordinate.  Thus inversion-fixed and paired classes have
stabilizers $2d$ and $d$, respectively; coefficient Frobenius is the
$p$-power map.  For $d=8$, it fuses the two odd pairs exactly when
$p\equiv\pm3\pmod8$.  On the tangent fiber the unipotent action is
$z\mapsto z+8u$, which vanishes exactly in characteristic two.
\end{proof}

\begin{proposition}[Other modular lifts]\label{prop:r9-other-modular}
The consecutive Lucas supports leave a line in characteristic five and
the quartic carrier above in characteristic seven; no other
characteristic has a nonzero coherent lift.  The former line is shallow
for every $q>5$.
\end{proposition}

\begin{proof}
Lucas' theorem and the consecutive-row overlap give the two supports.
In characteristic five, homogenize $t^5-t$ to include its simple root
at infinity and multiply by a linear factor $t-a$ with
$a\notin\F_5$.  The seven roots are distinct and the two required
Hankel coefficients vanish.  Such an $a$ exists for every $q>5$.
\end{proof}

\begin{proposition}[Degree-eight contained components]
\label{prop:r9-contained}
The assertion $\mathrm{CC}(8,1)$ holds.  The ordinary contained
component is the upper rank-two carrier.  The only modular lifts are
the characteristic-five point and the characteristic-seven
binary-quartic carrier of Proposition~\ref{prop:r9-components}; the
former is shallow for every field in the theorem range.  The marked
self-collision scheme is never the whole polar line.
\end{proposition}

\begin{proof}
Proposition~\ref{prop:contained-rank-two} gives the ordinary component,
and Proposition~\ref{prop:r9-other-modular} gives the complete Lucas
overlap and the characteristic-five witness.

For collision, remove the fixed divisor from the six-dimensional
Hankel series and let its moving degree be $d\leq7$.  If the associated
map were inseparable in characteristic $p$, its sections would lie in
the pullback of a series of degree $\lfloor d/p\rfloor$, of dimension
at most
\[
             \lfloor d/p\rfloor+1\leq4<6.
\]
Thus the moving series is separable.  Its ramification divisor has
degree at most $(5+1)(7-5)=12$, so collision is finite and supplies no
additional contained component.
\end{proof}

\begin{theorem}[Redundancy nine]\label{thm:r9}
For every prime power $q\geq53$, the
deep syndromes of $\PRS(q-8)$ are
exactly the persistent tangent and conjugate-secant families, of total
size $q(q+1)^2/2$.  Their orbit law is $T/T^8$ modulo inversion and
Frobenius; the tangent family splits exactly in characteristic two.
\end{theorem}

\begin{proof}
Proposition~\ref{prop:r9-budgets} handles every transverse
nonmodular point.  Proposition~\ref{prop:r9-contained} leaves the
persistent and modular supports.  Proposition
\ref{prop:r9-other-modular} eliminates characteristic five.  For
characteristic seven, the only powers below $343$ are $7$ and $49$,
both below the headline range; at every $q=7^m\geq343$,
Propositions~\ref{prop:r9-components} and \ref{prop:r9-deletion}
produce a witness.  The covering-radius theorem completes the
promotion to deep holes.
\end{proof}

Proposition~\ref{prop:r9-char7} is a carrier boundary, not part of the
headline classification: the $q=7$ count is not a deep-hole statement
for the inadmissible parameter $\PRS(-1)$, and the $q=49$ pass is not
an ambient $\PP^8$ census.
